\BourbakiExercisesSection

\begin{enumerate}
\def\labelenumi{\arabic{enumi}.}
\tightlist
\item
  设 E 是一个 A-代数，具有由两个元素 \(e_1, e_2\) 组成的基和一个单位元 \(e = \lambda e_1 + \mu e_2\) 。
\end{enumerate}

\begin{enumerate}
\def\labelenumi{(\alph{enumi})}
\item
  证明理想 \(a\) （ A 的由 \(\lambda\) 和 \(\mu\) 生成的理想）就是整个环 A（在相反的情形，注意 \(\mathrm{E} / \mathrm{aE} = \mathrm{E} \otimes_{\mathbf{A}} (\mathrm{A} / \mathrm{a}) = \{0\}\) 将与如下事实矛盾： \(\mathrm{A} / \mathrm{a} \neq \{0\}\) 且 E 是自由 A-模）。
\item
  若 \(\alpha, \beta\) 是 A 的两个元素，使得 \(\alpha\lambda + \beta\mu = 1\) ，证明元素 \(e\) 和 \(u = \beta e_1 - \alpha e_2\) 也构成 E 的一组基，并由此推出 E 是一个二次 A-代数。
\end{enumerate}

\begin{enumerate}
\def\labelenumi{\arabic{enumi}.}
\setcounter{enumi}{1}
\item
  证明 \(\mathbf{Z}\) 上的每个二次代数同构于类型为 \((0, n)\) 的代数（其中 \(n \in \mathbf{N}\) ），或类型为 \((m, 0)\) 的代数（其中 \(m\) 非平方数），或类型为 \((m, 1)\) 的代数（其中 \(m\) 是不具有形式 \(k(k - 1)\) 的整数）；此外，这些代数中任意两个都不同构。
\item
  \begin{enumerate}
  \def\labelenumii{(\alph{enumii})}
  \tightlist
  \item
    设 K 是一个交换域。证明：K 为 K 上类型为 \((\alpha, \beta, \gamma)\) 的四元数代数的中心的必要充分条件是以下情形之一成立：(1) \(\beta\gamma \neq 0\) ；(2) \(\gamma = 0\) , \(\beta \neq 0\) 且 \(\beta^{2} + 4\alpha \neq 0\) ；(3) \(\beta = 0\) , \(\alpha \neq 0\) 或 \(\gamma \neq 0\) ，且 \(2 \neq 0\) 在 K 中。
  \end{enumerate}
\end{enumerate}

\begin{enumerate}
\def\labelenumi{(\alph{enumi})}
\setcounter{enumi}{1}
\item
  设 K 是一个交换域，使得 \(2 \neq 0\) 在 K 中；证明 K 上类型为 \((1, \gamma)\) 的四元数代数同构于矩阵代数 \(\mathbf{M}_{2}(\mathbf{K})\) （考虑该代数由元素 \(\frac{1}{2}(1 + i)\) , \(\frac{1}{2}(1 - i)\) , \((j + k/2\gamma)\) , \(\frac{1}{2}(j - k)\) 组成的基）。
\item
  设 K 是一个交换域，使得 \(2 \neq 0\) 在 K 中，并设 E 为 K 上类型为 \((\alpha, \gamma)\) 的四元数代数。四元数 \(z \in E\) 称为纯的，如果 \(\overline{z} = -z\) （或等价地 \(T(z) = 0\) ）。证明：若 z 是纯四元数，则每个四元数 \(tzt^{-1}\) （其中 \(t \in E\) 可逆）也是纯的。
\item
  若 \(u, v\) 是任意两个四元数，则 \(uv - vu\) 是纯的。由此推出，若 \(u, v, w\) 是任意三个四元数，则
\end{enumerate}

\[
(u v - v u) ^ {2} w = w (u v - v u) ^ {2}.
\]

\begin{enumerate}
\def\labelenumi{(\alph{enumi})}
\setcounter{enumi}{4}
\tightlist
\item
  证明：若 E 中存在四元数 \(z \neq 0\) 使得 \(\mathrm{N}(z) = 0\) ，则也存在纯四元数 \(z' \neq 0\) 使得 \(\mathrm{N}(z') = 0\) 。（注意，在公式 (33)（第 5 号）的记号下，若 \(\alpha\) 不是 K 中的平方数，则 E 中存在 \(z \neq 0\) 使得 \(\mathrm{N}(z) = 0\) 这一事实等价于存在元素 \(y \in \mathrm{K} + \mathrm{Ki}\) 使得 \(\gamma = \mathrm{N}(y)\) 。）
\end{enumerate}

\begin{enumerate}
\def\labelenumi{\arabic{enumi}.}
\setcounter{enumi}{4}
\tightlist
\item
  在满足 2 = 0 的交换域 K 上，每个类型为 \((\alpha, 0, \gamma)\) 的四元数代数 E 都是交换的，且每个 \(x \in E\) 的平方都属于 K；因此，子代数 \(\mathrm{K}(x)\) （由 E 的元素 \(x \in K\) 生成）是 K 上类型为 \((\lambda, 0)\) 的二次代数，且 E 是 \(\mathrm{K}(x)\) 上的二次代数。
\end{enumerate}

证明集合 C （由 \(x \in E\) （其平方等于 K 中某个元素的平方）组成）是 E 的一个向量子空间，其维数等于 1、2 或 4。若 C 的维数为 1（此时 C = K），则 E 是一个域。若 C 的维数为 2，则存在一个包含于 E 中的二次代数 \(\mathrm{K}(x)\) ，它是一个域，并且 E 在 \(\mathrm{K}(x)\) 上有一组由两个元素 1 和 u 组成的基，满足 \(u^{2} = 0\) ； \(y \in E\) （满足 \(y^{2} = 0\) ）的集合是 K 上的一个 2 维理想 a，且 E/a 同构于 \(\mathrm{K}(x)\) 。最后，若 C 的维数为 4（从而等于 E），则 \(y \in E\) （满足 \(y^{2} = 0\) ）的集合 a 是一个 3 维理想，且 E/a 同构于 K；E 中存在一个基 \((1, e_{1}, e_{2}, e_{3})\) ，满足 \(e_{1}^{2} = e_{2}^{2} = e_{3}^{2} = 0\) , \(e_{1}e_{2} = e_{3}\) , \(e_{1}e_{3} = e_{2}e_{3} = 0\) ； \(Ke_{3} = b\) 是 E 中唯一的 1 维理想，b 是 a 的零化子，且 a/b 是 E/b 的两个 1 维理想的直和，这两个理想相互零化。

\begin{enumerate}
\def\labelenumi{\arabic{enumi}.}
\setcounter{enumi}{5}
\tightlist
\item
  设 \(\mathbf{K}\) 是一个满足 \(2 \neq 0\) 的交换域， \(\mathbf{E}\) 是 \(\mathbf{K}\) 上类型为 \((0, 0, \gamma)\) 的四元数代数。
\end{enumerate}

\begin{enumerate}
\def\labelenumi{(\alph{enumi})}
\item
  若 \(\gamma\) 不是 \(\mathbf{K}\) 中的平方数，则 \(\mathbf{E}\) 中不存在 \(\mathbf{K}\) 上维数为 1 的左（相应地，右）理想；集合 \(a\) （由 \(x \in \mathbf{E}\) （满足 \(x^2 = 0\) ）组成）是 \(\mathbf{K}\) 上维数为 2 的双边理想，且 \(\mathbf{E} / a\) 是一个域，即 \(\mathbf{K}\) 上的一个二次代数。
\item
  若 \(\gamma\) 是一个平方数 \(\neq 0\) （在 \(\mathbf{K}\) 中），则 \(\mathbf{E}\) 中存在一组基 \((e_i)_{1\leq i\leq 4}\) ，其乘法表如下：
\end{enumerate}

\(e_1\)

\(e_2\)

\(e_3\)

\(e_4\)

\(e_1\)

\(e_1\)

0

\(e_3\)

0

\(e_2\)

0

\(e_2\)

0

\(e_4\)

\(e_3\)

0

\(e_3\)

0

0

\(e_4\)

\(e_4\)

0

0

0

 \(x \in E\) （满足 \(x^{2} = 0\) ）的集合 a 是一个维数为 2 的双边理想，它是两个相互零化的双边理想 \(Ke_{3}\) 和 \(Ke_{4}\) 的直和；后者是 E 中仅有的维数为 1 的（左或右）理想。商代数 E/a 同构于两个与 K 同构的域的乘积。

\begin{enumerate}
\def\labelenumi{(\alph{enumi})}
\setcounter{enumi}{2}
\tightlist
\item
  若 \(\gamma = 0\) ，则集合 \(a\) （由 \(x \in E\) （满足 \(x^2 = 0\) ）组成）是一个维数为 3 的双边理想； \(Kk = b\) 是 \(E\) 中仅有的维数为 1 的（左或右）理想；它是一个双边理想，即 \(a\) 的左、右零化子。在商代数 \(E/b\) 中， \(a/b\) 是两个相互零化的 1 维双边理想的直和；最后， \(E/a\) 是一个与 \(K\) 同构的域。
\end{enumerate}

\begin{enumerate}
\def\labelenumi{\arabic{enumi}.}
\setcounter{enumi}{6}
\tightlist
\item
  设 K 是一个交换域，其中 \(2 \neq 0\) 且E为K上的代数，其基由四个元素1、i、j、k组成（1为单位元），乘法表如下：
\end{enumerate}

\[
i ^ {2} = j ^ {2} = k ^ {2} = 1, \quad i j = j i = k, \quad j k = k j = i, \quad k i = i k = j.
\]

证明 E 同构于四个与 K 同构的域的乘积（考虑由元素 \((1 + \varepsilon i)(1 + \varepsilon'j)\) （其中 \(\varepsilon\) 和 \(\varepsilon'\) 等于 1 或 -1）组成的 E 的基）。

代数 E 是（域 K 上的）两个 2 阶循环群的乘积群的代数。将此推广到 n 个 2 阶循环群的乘积群的代数（参见 § 4，第 1 号，例 2）。

\begin{enumerate}
\def\labelenumi{\arabic{enumi}.}
\setcounter{enumi}{7}
\item
  四元数群 \(\mathfrak{Q}\) （I，第 6 节，习题 4）同构于八个四元数 \(\pm 1\) , \(\pm i\) , \(\pm j\) , \(\pm k\) （在类型为 \((-1, -1)\) 的四元数代数（定义于满足 \(2 \neq 0\) 的域上）中）所成的乘法群。证明群 \(\mathfrak{Q}\) 在满足 \(2 \neq 0\) 的域 K 上的代数同构于四个与 K 同构的域之积以及 K 上类型为 \((+1, -1)\) 的四元数代数（若 \(c\) 是 \(\mathfrak{Q}\) 中对应于四元数 \(-1\) 的元素，则 \(\mathfrak{Q}\) 的元素可写成 \(e\) , \(i, j, k, c, ci, cj, ck\) ；考虑由元素 \(\frac{1}{2}(e + c)\) , \(\frac{1}{2}(e - c)\) , \(\frac{1}{2}(e + c)i\) , \(\frac{1}{2}(e - c)i\) , \(\frac{1}{2}(e + c)j\) , \(\frac{1}{2}(e - c)j\) , \(\frac{1}{2}(e + c)k\) , \(\frac{1}{2}(e - c)k\) 组成的 E 的基）。
\item
  \begin{enumerate}
  \def\labelenumii{(\alph{enumii})}
  \tightlist
  \item
    证明 8 阶二面体群 \(\mathbf{D}_4\) （I，第 6 节，习题 4）在满足 \(2 \neq 0\) 的域 K 上的代数 E 同构于四个与 K 同构的域的乘积以及 K 上的 2 阶矩阵代数。（若 \(a, b\) 是 I，第 6 节，习题 4 中引入的 \(\mathbf{D}_4\) 的两个生成元，则 \(\mathbf{D}_4\) 的元素可写成 \(a^i b^j\) ，其中 \(0 \leqslant i \leqslant 3\) , \(0 \leqslant j \leqslant 1\) ；考虑由四个元素 \(\frac{1}{2}(e + a^2)\) , \(\frac{1}{2}(e - a^2)\) , \(\frac{1}{2}(a + a^3)\) , \(\frac{1}{2}(a - a^3)\) 以及这四个元素右乘 \(b\) 所组成的 E 的基；利用习题 4。）由此推出，若 \(-1\) 是 K 中的平方数，则互不同构的群 \(\mathfrak{Q}\) 和 \(\mathbf{D}_4\) 在 K 上的代数是同构的。
  \end{enumerate}
\end{enumerate}

\begin{enumerate}
\def\labelenumi{(\alph{enumi})}
\setcounter{enumi}{1}
\tightlist
\item
  类似地证明：6 阶二面体群 \(D_{3}\) 在满足 \(2 \neq 0\) 和 \(3 \neq 0\) 的域 K 上的代数 F 同构于两个与 K 同构的域的乘积以及 K 上的 2 阶矩阵代数（这里考虑由元素 \(e + a + a^{2}\) , \(a + a^{2} - 2e\) , \(a - a^{2}\) 以及这三个元素右乘 b 所组成的 F 的基）。
\end{enumerate}

\begin{enumerate}
\def\labelenumi{\arabic{enumi}.}
\setcounter{enumi}{9}
\item
  设 G 是一个群，H 是 G 的一个正规子群。证明：若 a 是群代数 \(\mathbf{A}^{(\mathbf{G})}\) 的由元素 ts - s（其中 t 遍历 H，s 遍历 G）生成的双边理想，则群代数 \(\mathbf{A}^{(\mathbf{G}/\mathbf{H})}\) 同构于商代数 \(\mathbf{A}^{(\mathbf{G})}/\mathbf{a}\) 。
\item
  设 A 为交换环，S 为具有单位元 e 的交换幺半群。假设在 A 上给定一个以 S 为算子域的外合成律 \((s, x) \mapsto x^{s}\) ，使得对所有 \(s \in S\) ，映射 \(x \mapsto x^{s}\) 都是环 A 的自同构，且 \((x^{s})^{t} = x^{st}\) 对所有 \(x \in A\) 以及 S 中的 s 和 t 成立（这蕴含 e 是所考虑的外合成律的单位算子）。则在 A-模 \(\mathrm{A}^{(\mathrm{S})}\) 上定义了一个乘法的内部合成律，其典范基记为 \((b_{\mathrm{s}})\) ，由关系式
\end{enumerate}

\[
\left(\sum_ {s} x _ {s} b _ {s}\right) \left(\sum_ {s} y _ {s} b _ {s}\right) = \sum_ {s} \left(\sum_ {t u = s} a _ {t, u} x _ {t} y _ {u} ^ {s}\right) b _ {s}
\]

其中 \(a_{s,t}\) 属于 A。

证明这个合成律与 \(\mathbf{A}^{(\mathbf{S})}\) 上的加法一起在此集合上定义了一个环结构，只要 \(a_{s,t}\) 满足条件

\[
a _ {s, t} a _ {s t, u} = a _ {t u} ^ {s} a _ {s, t u}
\]

对于 S 中的所有 s、t、u。若还满足以下条件，则元素 \(b_{e}\) 是这个环的单位元

\[
a _ {e, s} = a _ {s, e} = 1
\]

对所有 \(s \in S\) 成立；此时 A 可等同于 \(A^{(S)}\) 的一个子环，并且若用 C 表示 A 的由在所有自同构 \(x \mapsto x^s\) 下不变的元素组成的子环，则 \(A^{(S)}\) 是一个 C-代数，称为环 A 与幺半群 S 关于因子系 \((a_{s,t})\) 的交叉积。若把因子系 \((a_{s,t})\) 换成 \(\left(\frac{c_s c_t^s}{c_{st}} a_{s,t}\right)\) （其中，对所有 \(s \in S\) , \(c_s\) 是 A 的可逆元素，且 \(c_e = 1\) ），则所得到的新交叉积同构于由因子系 \((a_{s,t})\) 所定义的交叉积。

特别地，若取 A 为环 C 上的二次代数，且 S 为 2 阶循环群，例如 \(\{e, s\}\) ，使得 \(x^s = \bar{x}\) 对所有 \(x \in \mathbf{A}\) 成立，证明 A 与 S 的每个交叉积都是 C 上的四元数代数。

\begin{enumerate}
\def\labelenumi{\arabic{enumi}.}
\setcounter{enumi}{11}
\tightlist
\item
  设 I 为非空集合。证明自由广群 M(I) 的元素可以与序列 \((a_{i})_{1 \leqslant i \leqslant n}\) （其中 n 为任意整数 \(\geqslant 1\) ）等同，其中 \(a_{i}\) 或者是 I 的元素，或者是特殊符号 P（表示“开括号”），并且这些序列满足以下条件：
\end{enumerate}

\begin{enumerate}
\def\labelenumi{(\roman{enumi})}
\item
  序列的长度 \(n\) 等于其中出现的符号 \(\mathbf{P}\) 的个数的两倍再加 1。
\item
  对所有 \(k \leqslant n - 1\) ，符号 \(\mathbf{P}\) （出现在子序列 \((a_i)_{1 \leqslant i \leqslant k}\) 中的）的个数 \(\geqslant k / 2\) 。
\end{enumerate}

（利用《集合论》I，附录，将 M(I) 的一个元素与每个序列 \((a_i)\) 相对应；例如，序列 PPPxyzPxy 对应于字 (((xy)z)(xy))。）

¶ 13. 设 A 是一个交换环，E 是一个 A-模。A-模 \(E_{n}\) 由下式归纳地定义：

\[
\mathbf {E} _ {1} = \mathbf {E}, \quad \mathbf {E} _ {n} = \bigoplus_ {p + q = n} \left(\mathbf {E} _ {p} \otimes_ {\mathbf {A}} \mathbf {E} _ {q}\right) \quad \text { for } n \geqslant 2.
\]

我们记 \(\mathbf{ME} = \bigoplus_{n=1}^{\infty} \mathbf{E}_n\) ，并在 ME 上按下述规则定义一个 A-代数结构：对 \(x_p \in \mathbf{E}_p\) 和 \(x_q \in \mathbf{E}_q\) ， \(x_p x_q = x_p \otimes x_q \in \mathbf{E}_{p+q}\) 。

\begin{enumerate}
\def\labelenumi{(\alph{enumi})}
\item
  证明：对每个 A-代数 B 和每个 A-线性映射 \(f: \mathbf{E} \to \mathbf{B}\) ，存在唯一的代数的 A-同态 \(\mathbf{ME} \to \mathbf{B}\) 延拓 \(f\) 。ME 称为 A-模 E 的自由代数。
\item
  整数 \(a(n)\) 由下列公式对 \(n\) 归纳地定义
\end{enumerate}

\[
a (1), \quad a (n) = \sum_ {p + q = n} a (p) a (q) \quad \text { if } n \geqslant 2.
\]

证明 \(E_{n}\) 同构于 \(a(n)\) 个同构于 \(E^{\otimes n}\) 的模的直和。

\begin{enumerate}
\def\labelenumi{(\alph{enumi})}
\setcounter{enumi}{2}
\tightlist
\item
  设 \(f(\mathbf{T})\) 是形式幂级数 \(\sum_{n=1}^{\infty} a(n)\mathbf{T}^n\) 。证明
\end{enumerate}

\[
f (\mathbf {T}) = \mathbf {T} + (f (\mathbf {T})) ^ {2}.
\]

*由此推出公式

\[
f (\mathrm{T}) = \frac {1 - \sqrt {1 - 4 \mathrm{T}}}{2} \quad \text { and } \quad a (n) = 2 ^ {n - 1} \frac {1 . 3 . 5 \dots (2 n - 1)}{n !}. *
\]

\begin{enumerate}
\def\labelenumi{(\alph{enumi})}
\setcounter{enumi}{3}
\tightlist
\item
  若 I 是任意集合且 \(E = A^{(I)}\) ，证明 \(\operatorname{Lib}_{\mathbf{A}}(\mathbf{I})\) 可与 ME 等同。存在典范同态 \(\mathbf{M}(\mathbf{I}) \xrightarrow{w} \mathbf{N}^{(\mathbf{I})} \xrightarrow{l} \mathbf{N}\) ，其复合是 \(\mathbf{M}(\mathbf{I})\) 中的长度； \(\omega(m)\) 称为 \(m \in \mathbf{M}(\mathbf{I})\) 的多重次数；映射 \(\omega\) 在 \(\operatorname{Lib}_{\mathbf{A}}(\mathbf{I})\) 上定义了一个类型为 \(N^{(I)}\) 的分次。证明在此分次下，集合 \((\operatorname{Lib}_{\mathbf{A}}(\mathbf{I}))_{\alpha}\) （由多重次数为 \(\alpha \in N^{(I)}\) 的齐次元素组成）是一个自由 A-模，其秩等于
\end{enumerate}

\[
a (\alpha) = a (n) \frac {n !}{\alpha !} = 2 ^ {n - 1} \frac {1 . 3 . 5 \dots (2 n - 1)}{\alpha !}
\]

（其中 \(\alpha! = \prod_{i \in I} (\alpha(i))!\) ，而 \(n = \sum_{i \in I} \alpha(i) = l(\alpha)\) 是 \(\alpha\) 的长度。

\begin{enumerate}
\def\labelenumi{\arabic{enumi}.}
\setcounter{enumi}{13}
\tightlist
\item
  \begin{enumerate}
  \def\labelenumii{(\alph{enumii})}
  \tightlist
  \item
    设 M 为一个幺半群，其合成律记为 \(\top\) ；进一步假设 M 被一个序关系 \(x \leqslant y\) 全序化，使得关系 \(x < y\) 与 \(x' \leqslant y'\) ，或者 \(x \leqslant y\) 与 \(x' < y'\) ，蕴含 \(x \top x' < y \top y'\) 。证明：若 A 是整环，则幺半群 M 在 A 上的代数没有零因子。
  \end{enumerate}
\end{enumerate}

\begin{enumerate}
\def\labelenumi{(\alph{enumi})}
\setcounter{enumi}{1}
\item
  由此推出，如果 \(\mathbf{A}\) 是整环，则每个多项式代数 \(\mathbf{A}[(\mathbf{X}_i)_{i\in I}]\) 是整环。
\item
  若 A 是整环，且 f 和 g 是 \(\mathrm{A}[(\mathrm{X}_{i})_{i\in\mathbf{I}}]\) 中的两个多项式，使得 fg 是齐次的且 \(\neq0\) ，证明 f 和 g 都是齐次的。特别地，环 \(\mathrm{A}[(\mathrm{X}_{i})_{i\in\mathbf{I}}]\) 的可逆元素正是环 A 的可逆元素。
\end{enumerate}

\begin{enumerate}
\def\labelenumi{\arabic{enumi}.}
\setcounter{enumi}{14}
\item
  设 A 为交换环，n 为整数 \(\geqslant2\) ，使得对所有 \(a\in A\) ，方程 nx=a 有解 \(x\in A\) 。设 m 为整数 \(\geqslant1\) ，u 为 A{[}X{]} 中形如 \(X^{mn}+f(X)\) 的多项式，其中 f 的次数 \(\leqslant mn-1\) ；证明存在多项式 \(v\in A[X]\) （形如 \(X^{m}+w\) ，其中 w 的次数 \(\leqslant m-1\) ），使得 \(u-v^{n}\) 为零或次数 \(<m(n-1)\) 。
\item
  设 A 是一个交换环， \(u = \sum_{k=0}^{m} a_k X^k\) 是环 A{[}X{]} 中的一个零因子。证明：若 \(v = \sum_{k=0}^{n} b_k X^k\) 是 \(\neq 0\) 的、A{[}X{]} 中次数为 \(n\) 的元素，使得 \(uv = 0\) ，则存在多项式 \(w \neq 0\) （次数 \(\leqslant n - 1\) ），使得 \(uw = 0\) 。（化归到 \(b_0 \neq 0\) 的情形；若 \(a_k v = 0\) （对 \(0 \leqslant k \leqslant m - 1\) ），证明可取 \(w = b_0\) ；若 \(a_k v = 0\) 对于
\end{enumerate}

\[
0 \leqslant k <   p \leqslant m - 1
\]

且 \(a_{p}v \neq 0\) ，证明 \(a_{p}b_{0} = 0\) ，因此可取 \(w = \sum_{k=0}^{n-1} a_{p}b_{k+1}\mathbf{X}^{k}\) 。由此推出存在元素 \(c \neq 0\) （属于 \(\mathbf{A}\) ），使得 \(cu = 0\) 。

\begin{enumerate}
\def\labelenumi{\arabic{enumi}.}
\setcounter{enumi}{16}
\item
  设 A 是一个交换环。证明在 A 上一元多项式的代数 A{[}X{]} 中，映射 \((u,v)\mapsto u(v)\) 是一个合成律，它关于 A{[}X{]} 中的加法和乘法是结合的且左分配的。若 A 是整环，证明关系 \(u(v)=0\) 蕴含 u=0 或 v 为常数，且若 \(u\neq0\) 而 v 的次数 \textgreater0 ，则 \(u(v)\) 的次数等于 u 与 v 的次数之积。
\item
  设 K 是一个交换域，K{[}X{]} 是 K 上单不定元的多项式环。证明环 K{[}X{]} 的每个自同构 s 保持 K 不变（习题 14(c)），并在 K 上诱导出该域的一个自同构；进一步证明必然有 \(s(\mathbf{X}) = a\mathbf{X} + b\) ，其中 \(a \neq 0\) （在 K 中）， \(b \in \mathbf{K}\) （参见习题 17）。反之，证明给定一个自同构 \(\sigma\) （ \(\mathbf{K}\) 的）和两个元素 \(a, b\) （ \(\mathbf{K}\) 的，满足 \(a \neq 0\) ），就定义了且仅定义了一个自同构 \(s\) （ \(\mathbf{K}[\mathbf{X}]\) 的），使得 \(s|_{\mathbf{K}} = \sigma\) 且 \(s(\mathbf{X}) = a\mathbf{X} + b\) 。若 \(G\) 是环 \(\mathbf{K}[\mathbf{X}]\) 的所有自同构构成的群， \(N\) 是 \(G\) 的由 \(\mathbf{K}[\mathbf{X}]\) 的 K-代数结构组成的子群，证明 \(N\) 是 \(G\) 的一个正规子群，且 \(G / N\) 同构于域 \(K\) 的自同构群；群 \(N\) 同构于定义在集合 \(K^* \times K\) 上、以合成律 \((\lambda, \mu)(\lambda', \mu') = (\lambda\lambda', \lambda' \mu + \mu')\) 定义的群。
\item
  证明8个不定元中的多项式恒等式
\end{enumerate}

\[
\begin{array}{r l} & (\mathrm{X} _ {1} ^ {2} + \mathrm{X} _ {2} ^ {2} + \mathrm{X} _ {3} ^ {2} + \mathrm{X} _ {4} ^ {2}) (\mathrm{Y} _ {1} ^ {2} + \mathrm{Y} _ {2} ^ {2} + \mathrm{Y} _ {3} ^ {2} + \mathrm{Y} _ {4} ^ {2}) \\ & \qquad = (\mathrm{X} _ {1} \mathrm{Y} _ {1} - \mathrm{X} _ {2} \mathrm{Y} _ {2} - \mathrm{X} _ {3} \mathrm{Y} _ {3} - \mathrm{X} _ {4} \mathrm{Y} _ {4}) ^ {2} \\ & \qquad + (\mathrm{X} _ {1} \mathrm{Y} _ {2} + \mathrm{X} _ {2} \mathrm{Y} _ {1} + \mathrm{X} _ {3} \mathrm{Y} _ {4} - \mathrm{X} _ {4} \mathrm{Y} _ {3}) ^ {2} \\ & \qquad + (\mathrm{X} _ {1} \mathrm{Y} _ {3} + \mathrm{X} _ {3} \mathrm{Y} _ {1} + \mathrm{X} _ {4} \mathrm{Y} _ {2} - \mathrm{X} _ {2} \mathrm{Y} _ {4}) ^ {2} \\ & \qquad + (\mathrm{X} _ {1} \mathrm{Y} _ {4} + \mathrm{X} _ {4} \mathrm{Y} _ {1} + \mathrm{X} _ {2} \mathrm{Y} _ {3} - \mathrm{X} _ {3} \mathrm{Y} _ {2}) ^ {2}. \end{array}
\]

（将第 5 号的公式 (32) 应用于一个适当的四元数代数。）

\begin{enumerate}
\def\labelenumi{\arabic{enumi}.}
\setcounter{enumi}{19}
\tightlist
\item
  证明不存在关于6个不定元的多项式恒等式
\end{enumerate}

\[
(\mathrm{X} _ {1} ^ {2} + \mathrm{X} _ {2} ^ {2} + \mathrm{X} _ {3} ^ {2}) (\mathrm{Y} _ {2} ^ {2} + \mathrm{Y} _ {2} ^ {2} + \mathrm{Y} _ {3} ^ {2}) = \mathrm{P} ^ {2} + \mathrm{Q} ^ {2} + \mathrm{R} ^ {2}
\]

其中 P、Q、R 是系数在 Z 中且关于不定元 \(X_{i}\) 和 \(Y_{i}\) 的三个多项式。（注意 15 = 3·5 不能写成形式 \(p^{2} + q^{2} + r^{2}\) ，其中 p、q、r 为整数。）

\begin{enumerate}
\def\labelenumi{\arabic{enumi}.}
\setcounter{enumi}{20}
\tightlist
\item
  设 S 为具有单位元 e 的乘法幺半群，满足第 10 号中的条件 (D)，并且进一步满足以下两个条件：(1) 关系 st = e 蕴含 s = t = e；(2) 对所有 \(s \in S\) ，有限序列 \((t_{i})_{1 \leq i \leq n}\) （由不同于 e 且满足 \(t_{1}t_{2} \ldots t_{n} = s\) 的元素组成）的项数 n 小于一个仅依赖于 s 的有限数 \(v(s)\) 。
\end{enumerate}

证明：元素 \(x = \sum_{s \in S} \alpha_s s\) （属于 \(S\) 在环 \(A\) 上的全代数）可逆的必要充分条件是 \(\alpha_e\) 在 \(A\) 中可逆（化归到 \(\alpha_e = 1\) 的情形，并利用恒等式

\[
e - z ^ {n + 1} = (e - z) (e + z + \dots + z ^ {n})).
\]

\begin{enumerate}
\def\labelenumi{\arabic{enumi}.}
\setcounter{enumi}{21}
\item
  若 K 是一个交换域，证明在形式幂级数环 \(K[[X_{1}, X_{2}, \ldots, X_{p}]]\) 中只有一个极大理想，即不可逆元素的集合。另一方面，证明在多项式环 \(K[X_{1}, \ldots, X_{p}]\) 中，对 \(p \geqslant 1\) 总存在若干个不同的极大理想。
\item
  设 K 为一个交换域；证明不存在形式幂级数 \(u(\mathbf{X}, \mathbf{Y}) \in \mathbf{K}[[\mathbf{X}, \mathbf{Y}]]\) 使得，对于某个整数 \(m > 0\) ,
\end{enumerate}

\[
(\mathrm{X} + \mathrm{Y}) u (\mathrm{X}, \mathrm{Y}) = \mathrm{X} ^ {m} \mathrm{Y} ^ {m}.
\]

\begin{enumerate}
\def\labelenumi{\arabic{enumi}.}
\setcounter{enumi}{23}
\item
  设 K 是一个交换域，并设 k 是一个使得 \(k \neq 0\) （在 K 中）的整数。证明对 K{[}{[}X{]}{]} 的每个常数项等于 1 的形式幂级数 u，存在形式幂级数 \(v \in K[[X]]\) ，使得 \(v^{k} = u\) 。
\item
  设 A 为一个环（不管交换与否），并设 \(\sigma\) 为 A 的一个自同态。在加法乘积群 \(\mathbf{E} = \mathbf{A}^{\mathbf{N}}\) 上定义一个内部合成律，写作 \((\alpha_{n})(\beta_{n}) = (\gamma_{n})\) ，其中 \(\gamma_{n} = \sum_{p + q = n}\alpha_{p}\sigma^{p}(\beta_{q})\) （约定 \(\sigma^0 (\xi) = \xi\) ）。
\end{enumerate}

\begin{enumerate}
\def\labelenumi{(\alph{enumi})}
\tightlist
\item
  证明这个合成律与 E 上的加法一起，在 E 上定义了一个环结构，并以序列 \((\alpha_{n})\) （其中 \(\alpha_{0}=1\) 和 \(\alpha_{n}=0\) ，对 \(n\geqslant1\) ）为单位元 e。把每个 \(\xi\in A\) 对应到元素 \((\alpha_{n})\in E\) （满足 \(\alpha_{0}=\xi,\alpha_{n}=0\) ，对 \(n\geqslant1\) ）的映射是 A 到一个
\end{enumerate}

E 的子环，并与之等同。于是我们记 \(\sum_{n=0}^{\infty} \alpha_n X^n\) 而不是

\[
u = \sum_ {n = 0} ^ {\infty} \alpha_ {n} X ^ {n} \neq 0.
\]

\((\alpha_{n})\) ，并且 \(\mathbf{X}^{p}\beta = \sigma^{p}(\beta)\mathbf{X}^{p}\) 对所有 \(\beta \in \mathbf{A}\) 成立；若 \(u = \sum_{n=0}^{\infty} \alpha_{n}\mathbf{X}^{n} \neq 0\) ，则最小整数 \(n\) （使得 \(\alpha_{n} \neq 0\) ）称为 \(u\) 的阶，记作 \(\omega(u)\) 。

\begin{enumerate}
\def\labelenumi{(\alph{enumi})}
\setcounter{enumi}{1}
\tightlist
\item
  若 A 是一个没有零因子的环，且 \(\sigma\) 是 A 到 A 的一个子环上的同构，证明 E 是一个无零因子的环，并且
\end{enumerate}

\[
\omega (u v) = \omega (u) + \omega (v)
\]

对 \(u \neq 0\) 和 \(v \neq 0\) 。

\begin{enumerate}
\def\labelenumi{(\alph{enumi})}
\setcounter{enumi}{2}
\item
  \(u = \sum_{n=0}^{\infty} \alpha_n X^n\) 可逆的必要充分条件是 \(\alpha_0\) 在 A 中可逆。
\item
  假设 A 是一个域，且 \(\sigma\) 是 A 的一个自同构。证明 E 具有一个左分式域（I，§ 9，习题 15），并且在这个域 F 中，每个元素 \(\neq0\) 都可以唯一地写成 \(uX^{-h}\) 的形式，其中 u 是 E 中阶为 0 的元素。
\end{enumerate}

¶ *26. (a) 设 A 和 B 是 R 的两个良序子集（它们必然是可数的；参见《一般拓扑学》IV，§ 2，习题 1）。证明 A + B 是良序的，并且对所有 \(c \in \mathbf{A} + \mathbf{B}\) ，只存在有限多个有序对 \((a, b)\) （满足 \(a \in \mathbf{A}\) , \(b \in \mathbf{B}\) 和 \(a + b = c\) ）（为证明 A + B 的非空子集有最小元素，考虑它在 R 中的最大下界）。

\begin{enumerate}
\def\labelenumi{(\alph{enumi})}
\setcounter{enumi}{1}
\tightlist
\item
  设 K 为一个交换域。在向量空间 \(K^{R}\) 中，考虑由元素 \((\alpha_{x})\) 组成的向量子空间 E，其中依赖于 \((\alpha_{x})\) 的、由 \(x \in \mathbf{R}\) （满足 \(\alpha_{x} \neq 0\) ）组成的集合是良序的。对 E 中任意两个元素 \((\alpha_{x})\) , \((\beta_{x})\) ，我们记 \((\alpha_{x})(\beta_{x}) = (\gamma_{x})\) ，其中 \(\gamma_{x} = \sum_{y+z=x} \alpha_{y} \beta_{z}\) （此和式由 (a) 保证有意义）。证明这个合成律与加法一起，在 E 上定义了一个域结构；E 的元素也记作 \(\sum_{t \in \mathbb{R}} \alpha_{t} X^{t}\) ，称为系数在 K 中、具有良序实指数的形式幂级数。*
\end{enumerate}

